\documentclass[11pt,twoside,openright]{book}

\usepackage[T1]{fontenc}
\usepackage[utf8]{inputenc}
\usepackage{amsmath,amssymb,amsthm,mathtools,bm}
\usepackage{mathpazo}
\usepackage{microtype}
\usepackage[inner=1.35in,outer=1.1in,top=1.25in,bottom=1.35in]{geometry}
\usepackage{booktabs}
\usepackage{fancyhdr}
\usepackage[hidelinks]{hyperref}
\usepackage{cleveref}

\setlength{\parindent}{1.4em}
\setlength{\parskip}{0pt}
\linespread{1.07}
\emergencystretch=3em

\pagestyle{fancy}
\fancyhf{}
\fancyhead[LE]{\small\thepage\quad\textsc{A Cinema of Individual Editions}}
\fancyhead[RO]{\small\nouppercase{\leftmark}\quad\thepage}
\renewcommand{\headrulewidth}{0.3pt}
\fancypagestyle{plain}{\fancyhf{}\fancyfoot[C]{\small\thepage}\renewcommand{\headrulewidth}{0pt}}

\theoremstyle{plain}
\newtheorem{theorem}{Theorem}[chapter]
\newtheorem{proposition}[theorem]{Proposition}
\newtheorem{lemma}[theorem]{Lemma}
\newtheorem{corollary}[theorem]{Corollary}
\theoremstyle{definition}
\newtheorem{definition}[theorem]{Definition}
\newtheorem{example}[theorem]{Example}
\theoremstyle{remark}
\newtheorem{remark}[theorem]{Remark}

\newcommand{\Capac}{\mathrm{Cap}}
\newcommand{\Rob}{\mathrm{Rob}}
\newcommand{\Inv}{\mathrm{Inv}}
\DeclareMathOperator{\loss}{loss}

\title{\Huge\bfseries A Cinema of Individual Editions\\[1.2ex]
\large Chapter Seventeen: Case Studies\\[0.4ex]
\large Working Draft}
\author{Flyxion\\[0.3ex]\small Independent Researcher}
\date{November 2026}

\begin{document}

\frontmatter
\maketitle
\tableofcontents
\mainmatter
\setcounter{chapter}{16}

\chapter{Case Studies}
\label{ch:cases}

Three studies carry the machinery of the book through complete examples. The first follows the engine scene of Chapter~3 through three local productions and a delayed recognition. The second follows a correction through the full ledger of a film that has two scenes, the lawn and a dishwashing scene. The third takes a measured passage of an existing film and follows the acoustic occupation record of Chapter~6 through a correction to its own reference. Each study states what the theory predicts, what it does not determine, and what an experiment would have to settle.

The dishwashing scene appears in no earlier chapter, and its pools are invented here for the purpose, with the invention stated where it is used. The other two studies use the scenes and data of the earlier chapters unchanged.

\section{Study One: The Engine Scene in Three Locales}\label{sec:engine17}

The engine scene has $126$ class tuples, of which $32$ are faithful to the script history, and four readings that a viewer may reach: diligence, negligence, disclosure, and concealment (Chapter~3). Release it in three locales, with local trace in the close shots and the handoff, so that the editions of different locales are different global classes (Chapter~11), and give each locale thirty viewers, drawn at random from its faithful editions. The audience is ninety, and the three locales together have $96$ faithful global editions, so the exact promise fits. An audience of forty in each locale would need $120$ and does not.

\textbf{Delayed recognition.} The faithful editions leave different numbers of readings open, and so do the viewers who hold them. \Cref{tab:eng17} gives the ninety viewers by the number of readings their edition leaves open at the first viewing. Twenty-one are settled from the start, and sixty-nine hold an edition that leaves at least two readings open. A later film that commits to the concealment, in the sense of Chapter~14, narrows the reading set of each of those sixty-nine, and leaves all ninety with the reading of concealment alone.

\begin{table}[ht]
\centering
\begin{tabular}{@{}lcccc@{}}
\toprule
Readings left open & 1 & 2 & 3 & 4\\
\midrule
Viewers (of 90) & 21 & 29 & 24 & 16\\
\bottomrule
\end{tabular}
\caption{Ninety viewers of the engine scene, thirty in each of three locales, by the number of readings their edition leaves open at the first viewing.}\label{tab:eng17}
\end{table}

\textbf{What two viewers can discuss.} The opening argument of the book is that a shared premiere still gives viewers something to talk about, even when their editions differ. The engine scene can say how much. Take the $32$ faithful editions as the editions of one locale, so that two viewers who see the same class see the same shot, and ask what each viewer is warranted in saying about the third bolt, the claim that it was left loose. An edition warrants a claim when every history it commits to satisfies the claim. Before any reveal film, $8$ of the $32$ editions show the loose third bolt directly, none warrants the claim by inference from other shots, and $24$ leave it open. Of the $496$ pairs of viewers, $28$ have both seen it, $192$ have one who has seen it and one who has not, and $276$ have neither. The evidence that two viewers can point to in conversation is thinner still: counting only the shots that bear on the story, the close-ups, the glance, and the handoff, $233$ of the $496$ pairs have no informative shot in common, $205$ have one, $55$ have two, and $3$ have three.

The reveal film changes what is shared and leaves the rest. A film that commits to the concealment settles the reading for all $32$ viewers, so every pair now shares the claim that someone skipped a bolt and concealed it, against $28$ pairs before. It does not settle which bolt. Only the $8$ viewers who were shown the third bolt loose still warrant the claim about it, and in $468$ of the $496$ pairs the viewers share the reading without both warranting the event. Three things are therefore different. Two viewers can share an event, in that both warrant the same claim about what happened. They can share evidence, in that both can point to the same shot. And they can share the experience of discovering it, which is the path of Chapter~14 and which the openness at first viewing in that chapter summarizes. The model counts the first two and does not measure the third. The script \texttt{verify\_ch17\_shared.py} reproduces these figures, and what viewers would actually say to one another remains an empirical question.

\textbf{A wrongful correction.} A viewer reports that the close shot of the third bolt shows it loose. By Chapter~7 the report is rejected when the departure is intended, as it is here, since the script history has the third bolt loose. Suppose it is instead upheld as an error and the class is removed. Each locale has $24$ faithful editions left, against an audience of thirty, and the shortfall is six in each locale. The class had been held by eight viewers in the first locale and seven in each of the others, and the free editions left over to move them to are two, one, and one. Eighteen of the ninety viewers cannot be given a distinct faithful edition in any repair. The editions that recover the concealment alone fall from eight to six.

\textbf{What the theory predicts.} It predicts the shortfall exactly, since the capacity of each locale is a count and the audience is another, and it predicts that the failure falls on the viewers who held the removed class. It predicts that the reveal film settles the reading for every viewer, since each faithful edition is consistent with the script history. It predicts that a rightful rejection of the report costs nothing.

\textbf{What it does not determine.} Whether a viewer in the sixty-nine finds the film open in the way the model counts, so that a settled viewer sees a different film. Whether three locales really differ at the close shots and the handoff, which was assumed. And what the viewers of a locale would do if told that eighteen of them share an edition.

\textbf{What would settle it.} A reading task with the faithful editions of one locale, with and without the reveal film beforehand, would show whether the reading sets of the model predict what viewers report. A comparison of the close shots across locales, by viewers who are asked to say whether the shots differ, would test the assumed local trace.

\section{Study Two: A Correction Through a Ledger}\label{sec:ledger17}

Take a film of two scenes, the lawn of production A and a dishwashing scene with four slots and, in order, three, three, two and three classes, two takes to a class. The pools of the dishwashing scene are invented for this study. At its third slot one class shows a plate clean before it has been washed, which is the same kind of continuity error that the agreement rule removes from the lawn. The editions of the film are pairs of a lawn edition and a dishwashing edition.

At the premiere the lawn has $494$ editions and the dishwashing scene $54$, so the film has $26{,}676$ and, with takes, $103{,}716{,}288$. After both corrections the lawn has $238$ and the dishwashing scene $27$, and the film has $6{,}426$ editions and $24{,}984{,}288$ with takes. The audience is five hundred viewers, each with a distinct pair, drawn at random from the premiere product.

The two corrections reach the ledger in sequence. \Cref{tab:ledger17} gives what each does. The lawn correction invalidates $263$ of the five hundred pairs, the dishwashing correction $228$, and both together $370$, which is the three quarters of the audience whose pair held a stale class in either scene.

\begin{table}[ht]
\centering
\begin{tabular}{@{}lrrrr@{}}
\toprule
Step & Pairs invalid & Radius & Cost & Records added\\
\midrule
Lawn correction & 263 & 2 & 310 & 263\\
Dishwashing correction & 228 & 2 & 237 & 228\\
\midrule
Total, in sequence & 370 viewers moved & & 547 & 491\\
One shot, both together & 370 & 3 & 539 & 370\\
Anticipated allocation & 0 & & 0 & 0\\
\bottomrule
\end{tabular}
\caption{A correction through a ledger of five hundred viewers of the two-scene film. Radius and cost are Hamming over the nine slots, survivors fixed, minimum total movement.}\label{tab:ledger17}
\end{table}

In sequence, $121$ of the $370$ viewers who move are moved twice, because a viewer placed by the first repair can land on a pair that the second removes. The ledger grows from $500$ records to $991$. A single repair of the two corrections together moves the same $370$ viewers for a slightly lower cost, $539$ against $547$, at a larger radius, three against two, and adds $370$ records, so the ledger has $870$. Allocating at the premiere from the pool that survives both corrections, which has $6{,}426$ pairs for five hundred viewers, moves no viewer and leaves five hundred records.

The archive has to keep every stale take that the ledger names. The five hundred viewers name all twelve stale takes of the lawn and both stale takes of the dishwashing scene, and the $370$ viewers whose editions held a stale class cannot be reconstructed if those takes are purged.

\textbf{What the theory predicts.} The counts of this section follow from the pools and the allocation, and the comparison of sequence and one shot follows from Chapter~10: the order in which corrections arrive changes the cost and the number of moves by viewers, and the final pool does not depend on it. The theory predicts that anticipating the corrections requires no later movement of any viewer here, because the surviving pool is nearly thirteen times the audience. It is not free: anticipation restricts the premiere catalogue from $26{,}676$ editions to $6{,}426$, which is the trade that the capacity of Chapter~9 describes.

\textbf{What it does not determine.} That the two corrections would arrive as separate events, which is a fact about a production's workflow. That a viewer who has been moved will be told, and what such a viewer's receipt should show, which Chapter~12 raised as a design choice and no count settles. That the dishwashing scene's continuity error has the structure assumed, since its pools are invented.

\textbf{What would settle it.} A production's own correction log, replayed against a ledger of its real audience, would give the numbers of the table for a real film. The comparison with anticipated allocation would then be a measure of what the production left on the table.

\section{Study Three: A Correction to the Record of a Measured Passage}\label{sec:passage17}

The acoustic occupation record of Chapter~6 is a measured account of one existing film, and the dialogue reference is one of its entries. It is the median of the level over the $1{,}499$ seconds that the subtitles define as dialogue, which is $-34.554$ LUFS, and everything else in the record is read against it: a second is occupied when it is more than ten decibels above it, a block is an occupied run of at least sixty seconds that merges gaps of up to five, and a recovery window is a run of at least five seconds at or below it or inaudible. With the recorded reference the film has nine blocks, $1{,}004$ seconds of them, and $105$ windows, $1{,}526$ seconds, and the passage of Chapter~6 has windows of $9$, $7$ and $13$ seconds and six blocks.

Suppose the reference is corrected. \Cref{tab:ref17} gives the recomputed record for a shift of the reference by amounts that correspond to real alternatives in the record: the proxy reference from $1{,}509$ proxy seconds, which is $0.44$ decibels lower, the reference taken from the centre channel, which is $2.33$ decibels lower, and shifts of the opposite sign of the same size.

\begin{table}[ht]
\centering
\begin{tabular}{@{}lrrp{2.5cm}rp{1.4cm}@{}}
\toprule
Reference shift & Blocks & Windows & Passage windows (s) & Slack (s) & Blocks in passage\\
\midrule
$-2.33$ dB (centre) & 9 & 85 & 5, 5 & 0 & 5\\
$-1.00$ dB & 9 & 100 & 9, 5, 9 & 8 & 6\\
$-0.44$ dB (proxy) & 9 & 100 & 9, 7, 13 & 14 & 6\\
$0$ (recorded) & 9 & 105 & 9, 7, 13 & 14 & 6\\
$+0.44$ dB & 7 & 98 & 9, 7, 13 & 14 & 5\\
$+1.00$ dB & 6 & 91 & 10, 12, 13 & 20 & 4\\
$+2.33$ dB & 2 & 94 & 6, 8, 10, 27, 6, 7, 9 & 38 & 2\\
\bottomrule
\end{tabular}
\caption{The occupation record of the measured film recomputed under shifts of the dialogue reference, with the rest of its definition unchanged.}\label{tab:ref17}
\end{table}

The record is sturdier in one direction than in the other. Lowering the reference by the proxy's $0.44$ decibels leaves the passage exactly as it was, with the same three windows and the same slack of fourteen seconds, and lowering it to the centre-channel value takes the slack to zero, since the passage then has two windows of five seconds with nothing to shorten. Raising it by the same $0.44$ decibels removes two of the nine blocks, one of them in the passage, and raising it by $2.33$ leaves two blocks. A half decibel is a small difference in a measurement of loudness and an entry in the record that decides whether a passage has a recovery pause worth preserving.

The reference is also sensitive to the alignment of the subtitles on which it rests. Shifting the dialogue seconds by minus three, minus one, plus one, plus three and plus ten seconds moves the median by $-0.45$, $-0.35$, $-0.65$, $-1.05$ and $-0.38$ decibels, downward in each of the five, and an offset of plus three seconds already changes the passage's windows to $9$, $5$ and $9$ with a slack of eight. A misaligned subtitle file therefore corrects the record in the direction that lowers the reference, and a lowered reference makes the film's blocks longer and its windows fewer.

\textbf{What the theory predicts.} The timing editions of Chapter~6 inherit the record. Slack bounds how far a cut can shorten a window before it stops qualifying, so a record whose passage slack goes from fourteen seconds to zero removes the room for any shortening, and a record whose slack goes to eight leaves a little. The corrected record changes the count of timing editions that stay within the window tolerance, though the counts of Chapter~6 for the slot handles do not change.

\textbf{What it does not determine.} Whether the recorded reference, the proxy, or the centre-channel value is the correct one, which depends on what a viewer takes to be the level of dialogue. Whether a window of five seconds at or below the reference is a pause that a viewer experiences as recovery, which the threshold of five seconds assumes and the record does not test. And whether two of the nine blocks that the upward shift removes were ever occupied in the sense of a listener.

\textbf{What would settle it.} A listening experiment on the passage, with the pauses of the three windows shortened by amounts within and beyond the slack, and with the blocks listened to at the recorded and the shifted reference, would show which of the records predicts what listeners report. The experiment would choose among the references, and so correct the record by a measurement. It is a different thing from the recomputation of the table. Changing the analytical reference changes the annotation of a fixed recording, whereas editing the sound tests what a listener experiences, and only the second can say which annotation is right.

\section{What the Studies Show and What They Assume}\label{sec:limits17}

The first study shows that a count settles whether a correction can be repaired and who is left without an edition, and that a wrongful correction costs a whole locale its promise. The second shows that a correction to two scenes costs about three quarters of a film's audience a move and leaves a ledger of twice the length, and that a pool that survives every correction would have cost nobody anything. The third shows that the record on which a timing edition rests is itself sensitive to a half decibel and to a few seconds of subtitle offset. None of the three is an experiment, and in each the last paragraph says what the experiment would be.

The assumptions are those of the earlier chapters, restated for each study where it matters: the Hamming dissimilarity on class tuples, the local trace of the locales, the invented pools of the dishwashing scene, the median reference and the thresholds of the occupation record. The studies were chosen to cover a delayed reading, a ledger, and a measurement, and a different choice of scene would have shown different numbers.

The last chapter of the book collects what remains open.

The scripts \texttt{lawnmodel.py} and \texttt{verify\_ch17.py}, with the data of Chapter~6, reproduce each figure of this chapter.

\begin{thebibliography}{99}

\bibitem{flyxion2026gf}
Flyxion.
\newblock \emph{Generative Film: Categories, Sheaves, and Types for Cinema}.
\newblock Unpublished monograph, October 2026.

\bibitem{flyxion2026ao}
Flyxion.
\newblock \emph{Acoustic Occupation: Salience, Recovery, and the Temporal Organization of Film Sound}.
\newblock Unpublished monograph, September 2026.

\end{thebibliography}

\end{document}
